\documentclass[10pt]{article} % use larger type; default would be 10pt

\usepackage{enumerate}
\usepackage{amsmath}
\usepackage{ulem}
\usepackage{hyperref}
\usepackage{circuitikz}
\input{.common.tex}

\title{Homework 5;
23.33/50.00 (46.66\%)
}
\begin{document}

\maketitle

\section*{Problem 1 \score{3.33}{10}}
\begin{enumerate}[(a)]
	\item NOT ok (SOP is wrong, POS is correct)
	\item NOT ok (what are essential minterms for $A'C$, $A'CD$ or $A'B'C$?)
	\item OK
\end{enumerate}
\section*{Problem 2 \score{10}{10}}
OK
\section*{Problem 3 \score{10}{10}}
OK, but why did you draw Carnaugh diagram two times?
\section*{Problem 4 \score{0}{10}}
\begin{enumerate}[(a)]
	\item NOT ok ($ACD'$ is not prime as it is included in $ad'$; note that don't-cares are
		treated as 1s, see p.~140);
	\item NOT ok ($b c'd$ is prime, but not mentioned);
\end{enumerate}
\section*{Problem 5 \score{0}{10}}
NOT ok (why $(a,b,c,d)=(0,1,0,0)$ yields 0? doesn't it correspond to number 2, which is prime?)
\end{document}
